\documentclass[12pt,a4paper]{article}
\usepackage[utf8]{inputenc}
\usepackage[spanish]{babel}
\usepackage{geometry}
\usepackage{listings}
\usepackage{xcolor}
\usepackage{parskip}
\usepackage{amsmath}
\usepackage{hyperref}

\geometry{margin=2.5cm}

\lstset{
  basicstyle=\ttfamily\small,
  keywordstyle=\color{blue}\bfseries,
  commentstyle=\color{gray},
  breaklines=true, frame=single, language=Java,
  numbers=left, numberstyle=\tiny\color{gray}
}

\title{\textbf{Algoritmo de Wang --- Demostrador de Teoremas v2}\\
  \large Documentación Técnica}
\author{Inteligencia Artificial}
\date{}

\begin{document}
\maketitle
\tableofcontents
\newpage

\section{Descripción del Proyecto}
Segunda versión mejorada del demostrador de Wang. Agrega UI enriquecida con salida a colores, selector de ejemplos predefinidos, soporte de más reglas de inferencia (Modus Ponens, Silogismo Hipotético, De Morgan, Distribución, Contrapositiva) y visualización del árbol de prueba.

\section{Reglas de Inferencia Adicionales}

\subsection{Modus Ponens}
\[ \frac{A, \quad A \rightarrow B}{B} \]

\subsection{Silogismo Hipotético}
\[ \frac{A \rightarrow B, \quad B \rightarrow C}{A \rightarrow C} \]

\subsection{Leyes de De Morgan}
\[ \neg(A \wedge B) \equiv \neg A \vee \neg B \]
\[ \neg(A \vee B) \equiv \neg A \wedge \neg B \]

\subsection{Distribución}
\[ A \wedge (B \vee C) \equiv (A \wedge B) \vee (A \wedge C) \]

\subsection{Contrapositiva}
\[ (A \rightarrow B) \equiv (\neg B \rightarrow \neg A) \]

\section{Implementación de Reglas}

\begin{lstlisting}
String aplicarRegla(String formula, String regla) {
    switch (regla) {
        case "MODUS_PONENS":
            // si A y A->B están en antecedente,
            // agregar B
            return aplicarModusPonens(formula);

        case "DE_MORGAN_NEG_AND":
            // !A&B  -->  !A | !B
            if (formula.matches("!\\((.+)&(.+)\\)")) {
                String[] p = extraerOperandos(formula);
                return "!" + p[0] + " | !" + p[1];
            }
            break;

        case "CONTRAPOSITIVE":
            // A->B  -->  !B->!A
            if (formula.contains("->")) {
                String[] p = formula.split("->");
                return "!" + p[1] + "->" + "!" + p[0];
            }
            break;
    }
    return formula;
}
\end{lstlisting}

\section{Selector de Ejemplos}

La v2 incluye ejemplos predefinidos seleccionables desde un \texttt{JComboBox}:

\begin{itemize}
  \item Modus Ponens: $\{P, P \rightarrow Q\} \vdash Q$
  \item Silogismo Hipotético: $\{P \rightarrow Q, Q \rightarrow R\} \vdash P \rightarrow R$
  \item De Morgan: $\{\neg(P \wedge Q)\} \vdash \neg P \vee \neg Q$
  \item Distribución: $\{P \wedge (Q \vee R)\} \vdash (P \wedge Q) \vee (P \wedge R)$
\end{itemize}

\section{Visualización Mejorada}

\begin{itemize}
  \item Cada regla aplicada se muestra en color distinto en el panel de salida.
  \item Verde: rama cerrada (válida).
  \item Rojo: rama abierta (posiblemente refutada).
  \item Azul: regla aplicada (paso intermedio).
  \item El árbol de derivación se muestra con indentación para visualizar la estructura.
\end{itemize}

\section{Diferencias con v1}

\begin{center}
\begin{tabular}{|l|c|c|}
\hline
\textbf{Característica} & \textbf{v1} & \textbf{v2} \\
\hline
Reglas implementadas & 5 & 10+ \\
Ejemplos predefinidos & No & Sí \\
Salida a colores & No & Sí \\
Árbol de derivación & No & Sí \\
Bicondicional & No & Parcial \\
\hline
\end{tabular}
\end{center}

\end{document}
